\section{Discussion}
\label{sec:discussion}

\textbf{Why correlation outperforms IG.} Three explanations fit.
\begin{itemize}
  \item IG asks which bins would change a few sampled pixels' component
        probabilities. Correlation asks which ion images follow a component
        across the whole section, which is closer to what mSCF1 rewards.
  \item Sampling explains little: four- to sixteen-fold more IG pixels gained at
        most 0.024 against a gap of 0.05--0.25.
  \item mSCF1 favours mask-like segmentations.
\end{itemize}
Correlation is therefore an ablation finding about this pipeline: the
label-free segmentation carries the signal, and gradient attribution through a
dense VAE adds nothing to it. We do not claim that IG is generally worse, or
that segment-then-correlate, which is established
practice~\citep{alexandrov2011efficient}, is new.

\textbf{What remains of the spatial hypothesis.} In a dense VAE, no
neighbourhood intervention beat shuffled, nonlocal context, and smoothing the
segmentation did not help either. At $3\times3$ scale, neighbours are
near-copies of the centre, so there is little for the model to use. Spatial
information helped as a ranking statistic and in S3PL, which needs neighbour
input. Even in S3PL, correct placement added only an inconsistent $+0.04$.
Spatially aware \emph{ranking}, not spatially aware \emph{encoding}, is what
the evidence supports, hence the title.

\textbf{Reconstruction is not ranking quality.} PCA-5 reconstructed GBM better
than the VAE yet ranked peaks about equally well. Earlier pilots showed the
same pattern: context models and the spatial penalty changed the latents
without improving reconstruction. Reconstruction rewards fitting all intensity
variance, while peak quality here depends on how well a low-dimensional
segmentation separates the classes, in line with
\citet{mazarura2025autoencoder}.

\textbf{The proposal's aim.} The aim was autonomous, bias-free peak detection
with a lightweight model. It is met only in part.
\begin{itemize}
  \item Fitting and ranking are label-free, but the number of segments and the
        peak budget come from annotations, and BIC could not replace the first.
  \item The 1D model was cheaper than S3PL in time and memory on the section
        measured, but it did not benefit from its spatial input.
\end{itemize}
S3PL's instability means single-run comparisons against it can mislead. We
report its GBM gap as unresolved and select no favourable seeds or
checkpoints.

\textbf{Limitations.}
\begin{itemize}
  \item \textbf{Statistical units.} GBM has four PDX lines, CAC one biopsy and
        RCC eight subjects. The $p$-values describe these sections, not a
        population.
  \item \textbf{Class-informed settings.} $K$ and the peak budgets are
        class-informed.
  \item \textbf{The metric.} mSCF1 measures mask agreement, not molecular
        identity.
  \item \textbf{RCC} depends partly on its unannotated region.
  \item \textbf{Post hoc tests.} Tests S, V, N and R and the S3PL investigation
        were chosen after earlier results. Each was predeclared before its own
        outputs, and all outcomes are reported.
  \item \textbf{Untested.} Graph-based and ion-image spatial models,
        cross-section transfer and other IG settings were not tested.
\end{itemize}
